\documentclass[10pt,a4paper]{amsart}
\usepackage[margin=19mm]{geometry}
\usepackage{amsmath,amssymb,mathtools}
\usepackage[scr=rsfs,cal=euler]{mathalfa}
\usepackage{xfrac}
\usepackage[unicode,hidelinks,bookmarks=false]{hyperref}
\newcommand{\ie}{i.e.\ }
\newcommand{\cf}{cf.\ }
\newcommand{\resp}{respectively\ }
\newcommand{\Subsection}{\subsection}
\providecommand{\nobreakdash}{\nobreak\hbox{-}}
\setlength{\emergencystretch}{2em}
\multlinegap0pt
\usepackage{bm}
\usepackage{bbm}
\usepackage{dsfont}
\usepackage{xspace}
\usepackage{cancel}
\usepackage{mathtools}
\usepackage{amsthm}
\makeatletter
\@ifundefined{theorem}{\newtheorem{theorem}{Theorem}}{}
\@ifundefined{lemma}{\newtheorem{lemma}[theorem]{Lemma}}{}
\makeatother
\title[Convergence Rates for Stochastic Proximal and Projection Est]{Convergence Rates for Stochastic Proximal and Projection Estimators}
\author{Diego Morales, Pedro P\'erez-Aros, Emilio Vilches}
\date{Source: arXiv:2602.06750v3 (CC BY 4.0)}
\begin{document}
\maketitle

\noindent\emph{Source and licence.} Convergence Rates for Stochastic Proximal and Projection Estimators. Version arXiv:2602.06750v3 (CC BY 4.0), distributed under the Creative Commons Attribution 4.0 International licence, as recorded in the arXiv licence metadata for this exact version and shown on the abstract page https://arxiv.org/abs/2602.06750v3. Full rights notice: licenses/arxiv-2602.06750-CC-BY-4.0.txt.\par\smallskip

\noindent\textit{Editorial scope.} Counted unit: the Theorem whose source label is \texttt{convergence} in the article (source0.tex lines 235--241, source label \texttt{convergence}), delivered together with the article's own complete original proof of it (source0.tex lines 242--318, one \texttt{proof} environment of 77 lines). No mathematics is written, repaired or re-derived: statement and proof are the article's own text line for line. Required context retained uncounted: co-area (lines 594--610); eq:corner-null (lines 596--598). Every definition, display and label of those lines is retained unchanged. Omitted from this package and not redistributed: 1--234, 319--593, 599--1058. Nothing inside a retained excerpt is omitted or altered: every retained body is delivered exactly as the article's lines are written, so no omission receipt of any kind accompanies this package. Citations: the retained excerpts carry no citation command at all, so no bibliography entry of the article is transcribed and no reference is redistributed. Reference adaptation: none was needed and none was made; every reference inside the delivered unit resolves inside it, so no unresolved reference or citation is produced. Printed slips kept literal: no printed word, symbol, hypothesis or number of the counted statement or of its proof was altered. Layout and class: the article's own class is \texttt{amsart} with packages including fontenc, inputenc, amssymb, mathtools, microtype, enumitem, subcaption, hyperref, dsfont, mathrsfs; the wrapper is the lane's pinned \texttt{amsart} editorial wrapper at 10pt on A4, which restates the theorem-like environments the retained text needs and adds the article's own macro definitions that the retained text needs (none), with extra wrapper packages (bm, bbm, dsfont, xspace, cancel, mathtools). The delivered numbering is the unit's own. Assets: no retained span contains a figure environment, a graphics-inclusion command, a PDF-page inclusion, a table or a TikZ picture; no image file is copied into this package, the package declares no asset dependency and no image file is redistributed with it. Licence: version arXiv:2602.06750v3 of this article is distributed under the Creative Commons Attribution 4.0 International licence, as recorded in the arXiv licence metadata for this exact version and shown on the abstract page https://arxiv.org/abs/2602.06750v3. A full rights notice is delivered at licenses/arxiv-2602.06750-CC-BY-4.0.txt. Characters: every retained excerpt is pure ASCII, so the delivered engine, XeLaTeX with its default Latin Modern text font, reports no missing character.
\begin{lemma}\label{co-area}
Let $D\subset \mathbb{R}^n$ be a closed convex set with nonempty interior, and for $R>0$ set $\Omega_R:=\operatorname{int}(D)\cap B_R$. Then $\Omega_R$ is a bounded Lipschitz domain. Moreover, for a.e. $R>0$,  
\begin{equation}\label{eq:corner-null}
\mathcal{H}^{n-1}\big(\partial D\cap \partial B_R\big)=0.    
\end{equation}
Let $F\in W^{1,1}_{\operatorname{loc}}(\operatorname{int}(D);\mathbb{R}^n)$. Then, for every $R>0$ such that \eqref{eq:corner-null} holds and $\int_{\Omega_{R}}(\vert F(x)\vert +\Vert \nabla F(x)\Vert)\,dx<\infty$, we have
\begin{equation*}
\int_{\Omega_R}\operatorname{div}F(x)\,dx
=
\int_{\partial D\cap B_R}F\cdot \nu_D\,d\mathcal{H}^{n-1}
+
\int_{\operatorname{int}(D)\cap \partial B_R} F\cdot \nu_{B_R}\,d\mathcal{H}^{n-1},
\end{equation*}
where $\nu_D$ is the outer unit normal to $D$, defined $\mathcal{H}^{n-1}$-a.e. on $\partial D$,
$\nu_{B_R}$ is the outer unit normal to $B_R$. Here, $F$ on $\partial\Omega_R$ denotes the \emph{interior} Sobolev trace of $F|_{\Omega_R}$ on $\partial\Omega_R$,
restricted to $\partial D\cap B_R$ and to $\operatorname{int}(D)\cap \partial B_R$.
\end{lemma}

\begin{theorem}\label{convergence}
   Let $f\colon \mathbb{R}^n \to \mathbb{R}\cup \{+\infty\}$ be proper, lower semicontinuous,  and $\rho$-weakly convex for some $\rho\geq 0$.  Assume that $\operatorname{dom}f$ has nonempty interior. Fix $x\in \mathbb{R}^n$ and $0<\lambda <1/\rho$, and set $\mu:=\frac{1}{\lambda}-\rho>0$. 
Then, for every $\delta>0$,
$$
\Vert m_{\delta}(x)-\operatorname{prox}_{\lambda f}(x)\Vert \leq \sqrt{\frac{n\delta}{\mu}}.
$$
\end{theorem}

\begin{proof} 
Fix $x\in \mathbb{R}^n$ and assume that $0<\lambda <1/\rho$. Consider 
\begin{equation*}
g(y):=f(y)+\frac{1}{2\lambda}\Vert x-y\Vert^2.
\end{equation*}
Since $f$ is $\rho$-weakly convex, it follows that $g$ is $\mu$-strongly convex and
\begin{equation}\label{growth}
    g(y)\geq g(\operatorname{prox}_{\lambda f}(x))+\frac{\mu}{2}\Vert y-\operatorname{prox}_{\lambda f}(x)\Vert^2 \textrm{ for all } y\in \mathbb{R}^n.
\end{equation}
Consider the vector field
$$
F(y):=(y-\operatorname{prox}_{\lambda f}(x))e^{-g(y)/\delta}.
$$
Now, let us show that  $\mathbb{E}_{\sigma_{\delta}}\left[\langle y-\operatorname{prox}_{\lambda f}(x),\nabla g(y)\rangle \right]\leq n\delta$.
 
Indeed, since $y\mapsto g(y)$ is proper, lsc, and convex, its domain is convex; in particular, $D:=\operatorname{dom}f=\operatorname{dom}g$ is convex with nonempty interior (by assumption). On $\operatorname{int}(D)$, every proper lsc convex function is locally Lipschitz, hence differentiable a.e. and belongs to $W^{1,1}_{\operatorname{loc}}(\operatorname{int}(D))$. Consequently, $\nabla g\in L^1_{\operatorname{loc}}(\operatorname{int}(D);\mathbb{R}^n)$, and for a.e. $y \in \operatorname{int}(D)$,
$$
\operatorname{div}F(y)=ne^{-g(y)/\delta}-\frac{1}{\delta}\langle y-\operatorname{prox}_{\lambda f}(x),\nabla g(y)\rangle e^{-g(y)/\delta}.
$$
Fix $R>0$ and set $\Omega_R:=\operatorname{int}(D) \cap B_R $. Since 
$$\Vert y-\operatorname{prox}_{\lambda f}(x)\Vert \leq R+\Vert \operatorname{prox}_{\lambda f}(x)\Vert =:C_R \textrm{ on } B_R
$$
and $e^{-g(y)/\delta}\leq e^{-\inf_{B_R}g/\delta}=:M_{R,\delta}<\infty$, it follows that for a.e. $y\in \Omega_{R}$,
$$
\vert \operatorname{div}F(y)\vert \leq nM_{R,\delta}+\frac{M_{R,\delta}C_{R}}{\delta}\Vert \nabla g(y)\Vert,
$$
Since $\nabla g\in L^1_{\operatorname{loc}}(\operatorname{int}(D);\mathbb{R}^n)$, we have $\nabla g\in L^1(\Omega_R;\mathbb{R}^n)$, and then $\operatorname{div}F\in L^1(\Omega_R)$. For $R>0$, apply the divergence theorem to $\Omega_R$:
\begin{equation*}
    \begin{aligned}
\int_{\Omega_R}\operatorname{div}F(y)dy&=\int_{\partial \Omega_R} F\cdot \nu\, d\mathcal{H}^{n-1}.
    \end{aligned}
\end{equation*}
Moreover, by Lemma \ref{co-area}, we can select $R_k\to \infty$ such that
\[
\int_{\partial \Omega_{R_k}} F \cdot \nu \, d\mathcal{H}^{n-1}
=
\underbrace{\int_{\partial D\cap {B}_{R_k}}
 F \cdot \nu_{D} \, d\mathcal{H}^{n-1}}_{=:I_1}+\underbrace{\int_{\partial {B}_{R_k}\cap D}
 F\cdot \nu_{{B}_{R_k}} \, d\mathcal{H}^{n-1}}_{=:I_2}.
\]
Now, on the one hand, since $D$ is convex and $\operatorname{prox}_{\lambda f}(x)\in D$, for $\mathcal{H}^{n-1}$-a.e. $y\in \partial D$ one has
$$
\langle y-\operatorname{prox}_{\lambda f}(x),\nu_D(y)\rangle \geq 0 \quad \Rightarrow \quad   F(y)\cdot \nu_D(y) \geq 0.
$$
Hence the $\partial D$-term  $I_1$ is nonnegative.  On the other hand, using \eqref{growth}, the term $I_2$  is bounded in absolute value by 
$$
CR_k^n \exp\left(-\frac{\mu}{2\delta}(R_k-\Vert \operatorname{prox}_{\lambda f}(x) \Vert)^2\right)\to 0 \textrm{ as } k\to \infty,
$$
hence the $\partial {B}_{R_k}$-term $I_2$ vanishes as $k\to \infty$. Therefore,
\begin{equation*}
\begin{aligned}
    0&\leq \limsup_{k\to \infty} \int_{ B_{R_k}\cap \partial D }  F\cdot \nu_D\, d\mathcal{H}^{n-1}\\
    &=\limsup_{k\to \infty}\Big(\int_{\partial D\cap  {B}_{R_k}  }  F\cdot \nu_D\, d\mathcal{H}^{n-1}+\int_{\partial {B}_{R_k}\cap D}  F\cdot \nu_{{B}_{R_k}}\, d\mathcal{H}^{n-1}\Big)\\
    &=\limsup_{k\to \infty}\int_{\partial \Omega_{R_k}} F\cdot \nu\, d\mathcal{H}^{n-1}\\
    &=\limsup_{k\to \infty}\int_{\Omega_{R_k}} \operatorname{div}F(y)dy\\
    &=\int_D \operatorname{div}F(y)dy\\
    &=n\int_D e^{-g(y)/\delta}dy-\frac{1}{\delta}\int_D \langle y-\operatorname{prox}_{\lambda f}(x),\nabla g(y)\rangle e^{-g(y)/\delta}dy.
\end{aligned}
\end{equation*}
Finally, dividing by $\int_{D}e^{-g(y)/\delta}dy$ and rearranging, we obtain the estimate. 
Next, by strong convexity and almost-everywhere differentiability, we obtain
$$
\langle \nabla g(y),y-\operatorname{prox}_{\lambda f}(x)\rangle \geq \mu \Vert y-\operatorname{prox}_{\lambda f}(x)\Vert^2.
$$
Taking expectation yields
$$
\mu \mathbb{E}_{\sigma_{\delta}}\left[\Vert y-\operatorname{prox}_{\lambda f}(x)\Vert^2\right]\leq\mathbb{E}_{\sigma_{\delta}}\left[\langle y-\operatorname{prox}_{\lambda f}(x),\nabla g(y)\rangle \right]=n\delta,
$$
hence $\mathbb{E}_{\sigma_{\delta}}\left[\Vert y-\operatorname{prox}_{\lambda f}(x)\Vert^2\right]\leq n\delta /\mu$. Finally, by Jensen's inequality, 
\begin{equation*}
    \begin{aligned}
  \Vert m_{\delta}(x)-\operatorname{prox}_{\lambda f}(x)\Vert^2&\leq \mathbb{E}_{\sigma_{\delta}}\left[\Vert y-\operatorname{prox}_{\lambda f}(x)\Vert\right]^2\\
  &\leq  \mathbb{E}_{\sigma_{\delta}}\left[\Vert y-\operatorname{prox}_{\lambda f}(x)\Vert^2\right]\\
  &\leq \frac{n\delta}{\mu}.    
    \end{aligned}
\end{equation*}
\end{proof}

\end{document}
